### Basic objects

Fix real parameters (t,\alpha\in(0,\infty)). Define the integer sequence
[
a_n:=\big\lfloor t,\alpha^n\big\rfloor,\qquad n\in\mathbb N.
]

**Convention for (\mathbb N):** (\mathbb N={1,2,3,\dots}).
(Starting at (n=0) instead changes only finitely many initial terms, which does not affect “all sufficiently large integers” questions.)

Define the **subset-sum set**
[
\Sigma(a):=\left{\sum_{j=1}^k a_{n_j}\ :\ k\ge 0,\ n_1,\dots,n_k\in\mathbb N\ \text{pairwise distinct}\right}\subseteq\mathbb Z_{\ge 0}.
]
(Empty sum (k=0) equals (0).)

### Completeness (as in additive number theory)

Call ((a_n)) **complete** (more precisely: *cofinitely complete*) if
[
\exists N_0\in\mathbb N\ \ \forall m\in\mathbb N\ (m\ge N_0\implies m\in \Sigma(a)).
]
Equivalently, (\mathbb N\setminus \Sigma(a)) is finite.

### Ambiguities / misstatements in the pasted problem text

The pasted text says “sum of **distinct integers** of the form (\lfloor t\alpha^n\rfloor).” There are two standard interpretations:

* **(Index-distinct / standard in “distinct terms” problems):** you may use each *term* (a_n) at most once, even if (a_n=a_{n'}) for (n\neq n'). This is exactly the (\Sigma(a)) defined above.
* **(Value-distinct / literal reading):** you may use each *value* (v\in{a_n:n\in\mathbb N}) at most once, so duplicates would be unusable multiple times. This changes the problem in some parameter regimes (duplicates can matter).

The Erdős-problems discussion notes that the “distinct integers” phrasing is misleading compared to the intended “distinct terms” interpretation in the classical literature on “complete sequences.”

### Minimal corrected “statement” to prove/disprove

As written, the problem is a **classification question** (“for what values…?”), not a proposition with a truth value. The only propositional claim present is the conjectural sentence:

> “It seems likely that the sequence is complete for all (t>0) and all (1<\alpha<\frac{1+\sqrt5}{2}).”

So the **minimal corrected statement** consistent with standard conventions is:

> **Conjecture (C).** For every (t>0) and every (\alpha\in(1,\varphi)) (where (\varphi=\frac{1+\sqrt5}{2})), the sequence (a_n=\lfloor t\alpha^n\rfloor) is complete (in the index-distinct sense).

My job below is to try to prove or disprove Conjecture (C). If I cannot, I must output **UNRESOLVED** with a crisp first gap and rigorous partial results.

---
